\documentclass[aspectratio=169]{beamer}

\usepackage[ngerman]{babel}
\usepackage{amsmath, amssymb}
\usepackage{graphicx}

\usetheme{Madrid}
\usecolortheme{beaver}

\title{Neovim Beamer-Integrationstest}
\subtitle{Folien, Overlays und Mathe}
\author{Test Benutzer}
\date{\today}

\begin{document}

\begin{frame}
    \titlepage
\end{frame}

\begin{frame}{Inhalt}
    \tableofcontents
\end{frame}

\section{Grundlagen}

\begin{frame}{Aufzählung mit Overlays}
    \begin{itemize}
        \item<1-> Erster Punkt
        \item<2-> Zweiter Punkt erscheint später
        \item<3-> Dritter Punkt
    \end{itemize}
    \only<3>{\alert{Alle Punkte sind sichtbar.}}
\end{frame}

\begin{frame}{Blöcke}
    \begin{block}{Normaler Block}
        Ein gewöhnlicher Hinweis.
    \end{block}
    \begin{alertblock}{Warnung}
        Ein hervorgehobener Hinweis.
    \end{alertblock}
    \begin{exampleblock}{Beispiel}
        Ein Beispiel.
    \end{exampleblock}
\end{frame}

\section{Mathematik}

\begin{frame}{Formeln}
    \begin{theorem}[Pythagoras]
        Für ein rechtwinkliges Dreieck gilt $a^2 + b^2 = c^2$.
    \end{theorem}
    \begin{equation}
        \label{eq:euler}
        e^{i\pi} + 1 = 0
    \end{equation}
    Siehe Gleichung~\eqref{eq:euler}.
\end{frame}

\begin{frame}{Zwei Spalten}
    \begin{columns}
        \column{0.5\textwidth}
        Linke Spalte mit Text.
        \column{0.5\textwidth}
        \[
            \int_0^1 x^2 \, dx = \frac{1}{3}
        \]
    \end{columns}
\end{frame}

\end{document}
